\chapter{Patterns: state machines for UI}
\label{ch:uistate}

\begin{aside}
A button's state: idle, hover, pressed, disabled. A login
form's state: empty, typing, validating, error, success. State
machines model these explicitly.
\end{aside}

\section{Why explicit state}

A boolean is rarely sufficient. ``is\_loading'' tells you one
thing but not what it means in context. A state machine names
each combination.

\section{Encoding}

\begin{lstlisting}
enum class LoginState {
    Empty, Typing, Validating, Error, Success
};
auto state = signal(LoginState::Empty);
\end{lstlisting}

\section{Transitions}

Each transition is a method or function:

\begin{lstlisting}
void on_submit() {
    if (state == Typing) state = Validating;
}
\end{lstlisting}

Transitions from invalid states are no-ops (or asserts).

\section{Side effects}

A transition can fire side effects (network request, log
event, animation). Keep side effects in transitions, not
in renders.

\section{Visualisation}

For complex machines, draw the diagram. Check it in. Keep
in sync.

\section{Statechart libraries}

For very complex flows: boost.sml, xstate-cpp. Declarative,
testable. \maya{} doesn't bundle one; pick when complexity
demands.

\section{Recap}

\begin{recap}
\begin{itemize}[leftmargin=*]
  \item Enum state, signal-wrapped.
  \item Transitions are methods.
  \item Side effects in transitions, not renders.
  \item Diagram complex machines.
\end{itemize}
\end{recap}

\section*{Workshop}

\begin{enumerate}[leftmargin=*]
  \item Convert one boolean flag to a state machine.
  \item Draw the diagram.
  \item Test each transition.
\end{enumerate}
